
\subsection*{1.6}
\phantomsection
\label{sga3:VIIB:1.6}
\addcontentsline{toc}{subsection}{\SGARefLink{sga3:VIIB:1.6}{1.6}. Étale formal varieties}

To complete these generalities on formal varieties, it remains to
define briefly the formal varieties that are étale over \(k\).
\footnote{\label{sga3:VIIB:nde:70}\textit{Editors' note.} The original continued as follows:
``A formal variety \(X\) over \(k\) is said to be étale if the
diagonal morphism
\(\Delta_X:X\to X\times X\) is a local isomorphism; that is, if
\(\Delta_X\) induces an isomorphism from
\(\mathcal O_{X\times X,\Delta_X(x)}\) onto
\(\mathcal O_{X,x}\) for every point \(x\) of \(X\). It follows
readily from SGA~I that this formulation is equivalent to each of the
following two: the formal variety \(X\) is topologically flat and,
for every point \(x\in X\) projecting to
\(s\in\operatorname{Spf}(k)\),
\(\mathcal O_{X,x}\widehat\otimes_k\kappa(s)\) is a finite separable
extension of the residue field \(\kappa(s)\); or, if \(A\) denotes the
affine algebra of \(X\), the local components \textup{(\SGARefLink{sga3:VIIB:0.1}{0.1})} of
\(A\widehat\otimes_k(k/\mathfrak l)\) are finite étale algebras over
\(k/\mathfrak l\), for every open ideal \(\mathfrak l\) of \(k\).''
In what follows, the editors corrected the omission of the flatness
hypothesis from the first condition above and supplied details for
the equivalence of the conditions.}

\medskip
\noindent\textbf{Recollections \SGARefLink{sga3:VIIB:1.6.A}{1.6.A}.}
\phantomsection
\label{sga3:VIIB:1.6.A}
\textit{(i)} Recall first \textup{(cf.\ EGA \(0_{\mathrm{IV}}\),
19.10.2)} that a topological \(k\)-algebra \(A\) is said to be
\emph{formally étale over \(k\)}---for the given topologies on \(k\)
and \(A\), respectively for the discrete topologies---if, for every
discrete topological \(k\)-algebra \(C\) \textup{(not necessarily
Artinian)} and every nilpotent ideal \(J\) of \(C\), every continuous
morphism of \(k\)-algebras
\[
  A\longrightarrow C/J
\]
lifts uniquely to a continuous morphism \(A\to C\), where \(A\) is
endowed with the given topology, respectively with the discrete
topology.

This property is plainly preserved under base change. Thus, for
every morphism \(k\to k'\) of pseudocompact rings,
\[
  A\widehat\otimes_k k'
\]
is formally étale over \(k'\). It is also enough to verify the
lifting condition for square-zero ideals \(J\); cf.\ EGA
\(\mathrm{IV}_4\), 17.1.2~(ii).

The algebra \(A\) is said to be \emph{étale over \(k\)} if it is
formally étale over \(k\) for the discrete topologies and is, in
addition, a \(k\)-algebra of finite presentation; cf.\ EGA
\(\mathrm{IV}_4\), 17.3.2~(ii). In what follows, when \(k\) is a
pseudocompact ring and \(A\) a profinite \(k\)-algebra, ``formally
étale'' refers to the given topologies unless the contrary is stated.

\smallskip
\noindent\textit{(ii)} Recall also that if \(F\in k[T]\) is a
separable monic polynomial of degree \(d\geq1\), meaning that the
ideal generated by \(F\) and its derivative \(F'\) is all of
\(k[T]\), then
\[
  B=k[T]/(F),
\]
which is free of rank \(d\) over \(k\) and is endowed with the product
topology, is formally étale over \(k\).

Indeed, let \(C\) be a discrete \(k\)-algebra, so that the kernel
\(\mathfrak l\) of \(k\to C\) is an open ideal of \(k\); let \(J\)
be a square-zero ideal of \(C\); and let
\(\phi:B\to C/J\) be a continuous morphism of \(k\)-algebras. Since
\(B\) is a free \(k\)-module of finite rank,
\(\mathfrak lB\) is an open ideal of \(B\), and hence every lift
\(\Phi:B\to C\) of \(\phi\) is automatically continuous.

Let \(t\) be the image of \(T\) in \(B\), and choose any lift
\(u_0\in C\) of \(\phi(t)\). Then \(F(u_0)\in J\), since
\(\phi(t)\) is a root of \(F\). There exist \(G,H\in k[T]\) such that
\[
  GF+HF'=1.
\]
Consequently,
\[
  H(u_0)F'(u_0)=1-G(u_0)F(u_0).
\]
The right-hand side is invertible because \(F(u_0)\) is square-zero,
and therefore \(F'(u_0)\) is invertible. We seek \(h\in J\) such that
\(u=u_0+h\) is a root of \(F\). This is equivalent to
\[
  0=F(u)=F(u_0)+F'(u_0)h,
\]
whose unique solution is
\[
  h=-F'(u_0)^{-1}F(u_0)\in J.
\]
The same proof, with no continuity assumptions on the morphisms
\(k\to C\), \(\phi\), and \(\Phi\), also shows that \(B\) is an
étale \(k\)-algebra.

\smallskip
\noindent\textit{(iii)} Finally, if
\[
  A=A_1\times\cdots\times A_n
\]
is a finite product, then \(A\) is formally étale over \(k\) if and
only if every \(A_i\) is.
It is enough to prove this for \(n=2\). Let \(e=1_{A_1}\) be the
idempotent such that \(A_1=Ae\) and \(A_2=A(1-e)\), and suppose that
a continuous morphism
\[
  \phi:A\longrightarrow C/J
\]
is given, where \(J\) is square-zero. The polynomial
\(F=X^2-X\) is separable, since
\[
  F'=2X-1
  \qquad\text{and}\qquad
  (F')^2-4F=1.
\]
Therefore the idempotent \(\phi(e)\) of \(C/J\) lifts uniquely to an
idempotent \(f\) of \(C\). Hence
\[
  C=Cf\oplus C(1-f),
\]
and giving a lift \(\Phi\) is equivalent to giving two morphisms
\[
  \Phi_1:A_1\longrightarrow Cf,
  \qquad
  \Phi_2:A_2\longrightarrow C(1-f),
\]
lifting the respective restrictions of \(\phi\). The same argument
shows that if \(e\) is an idempotent of \(k\) and \(A=Ae\), then
\(A\) is formally étale over \(k\) if and only if it is formally
étale over the local component \(k_e\), identified with \(ke\).

Now let \(X\) be a formal \(k\)-variety with affine algebra \(A\).
If \(x\) is a point of \(X\), that is, an open maximal ideal
\(\mathfrak m\) of \(A\), and \(s\) is its image in
\(\operatorname{Spf}(k)\), we denote by \(k_{\mathfrak m}\), or
\(k_s\), the local component of \(k\) corresponding to \(s\).

\medskip
\noindent\textbf{Definition 1.6.B.}
\phantomsection
\label{sga3:VIIB:1.6.B}
\textit{The following conditions are equivalent:}
\begin{enumerate}
\item[\textup{(a)}] \(A\) is formally étale over \(k\).
\item[\textup{(b)}] For every \(\mathfrak m\in\Upsilon(A)\),
  \(A_{\mathfrak m}\) is formally étale over \(k\), or equivalently
  over \(k_{\mathfrak m}\).
\item[\textup{(c)}] For every open ideal \(\mathfrak l\) of \(k\),
  \[
    A\widehat\otimes_k(k/\mathfrak l)
    =
    A/A\mathfrak l
  \]
  is formally étale over \(k/\mathfrak l\).
\item[\textup{(d)}] For every \(\mathfrak m\in\Upsilon(A)\) and every
  open ideal \(\mathfrak l\) of \(k_{\mathfrak m}\),
  \(A_{\mathfrak m}/A_{\mathfrak m}\mathfrak l\) is formally étale
  over \(k_{\mathfrak m}/\mathfrak l\).
\end{enumerate}

We say that \(X\) is \emph{étale over \(k\)} if these conditions
hold, and write
\[
  \operatorname{Vaf}^{\mathrm{\acute{e}t}}/k
\]
for the full subcategory of \(\operatorname{Vaf}/k\) consisting of
formal varieties that are étale over \(k\).

Indeed, suppose that
\[
  \phi:A\longrightarrow C/J
\]
is a continuous morphism of \(k\)-algebras, where \(C\) is a
discrete \(k\)-algebra and \(J\) is square-zero. Then
\(I=\operatorname{Ker}(\phi)\) is an open ideal of \(A\), so \(A/I\)
is Artinian. Thus \(I\) is contained in only finitely many open
maximal ideals
\[
  \mathfrak m_1,\ldots,\mathfrak m_r.
\]
It therefore contains the product of the components
\(A_{\mathfrak m}\) for
\(\mathfrak m\notin\{\mathfrak m_1,\ldots,\mathfrak m_r\}\). This
product is \(A(1-e)\), where \(e\) is the idempotent of \(A\) for
which
\[
  Ae=\prod_{i=1}^{r}A_{\mathfrak m_i}.
\]
Hence \(\phi(e)=1_{C/J}\), and giving a continuous lift of \(\phi\)
is equivalent to giving a continuous lift of the induced morphism
\[
  Ae\simeq A/A(1-e)\longrightarrow C/J.
\]

On the other hand,
\[
  A\simeq\varprojlim_{\mathfrak l}A/A\mathfrak l.
\]
The equivalence of the stated conditions follows readily from these
observations and the preceding recollections.

\medskip
\noindent\textbf{Definitions 1.6.C.}
\phantomsection
\label{sga3:VIIB:1.6.C}
Put
\[
  \kappa(k)
  =
  \prod_{s\in\operatorname{Spf}(k)}\kappa(s),
\]
endowed with the product topology. Thus the formal variety
\(\operatorname{Spf}(\kappa(k))\) is the direct sum of the schemes
\(\operatorname{Spec}\kappa(s)\), for
\(s\in\operatorname{Spf}(k)\). We write \(S_{\kappa(k)}\) for the
scheme that is the direct sum of the same
\(\operatorname{Spec}\kappa(s)\).

For every formal variety \(X\) over \(k\), put
\[
  X_\kappa
  =
  X\widehat\otimes_k\kappa(k).
\]
This is the formal variety over \(\kappa(k)\) obtained by base
change. It has the same points as \(X\), and if
\(x\in X\) projects to \(s\in\operatorname{Spf}(k)\), then
\[
  \mathcal O_{X_\kappa,x}
  =
  \mathcal O_{X,x}\widehat\otimes_k\kappa(s).
\]
The base-change functor
\[
  \operatorname{Vaf}/k
  \longrightarrow
  \operatorname{Vaf}/\kappa(k)
\]
sends \(\operatorname{Vaf}^{\mathrm{\acute{e}t}}/k\) into
\(\operatorname{Vaf}^{\mathrm{\acute{e}t}}/\kappa(k)\), by \SGARefLink{sga3:VIIB:1.6.A}{1.6.A}~\SGARefLink{sga3:vii:paragraph:5-1:equation:i}{(i)}.
